% Generated from locale/tr/content/first-order-logic/tableaux/rules-and-proofs.tex; do not hand-edit.


\iftag{FOL}
      {\olfileid[tr]{fol}{tab}{rul}}
      {\olfileid[tr]{pl}{tab}{rul}}

\olsection{Kurallar ve \usetoken{P}{tableau}}

!!^a{tableau}, !!a{sentence}'nin
\iftag{FOL}{!!a{structure}}{bir değerleme} içinde doğru ya da yanlış
olabileceği yolların sistematik bir incelemesidir. Bir tableau'nun yapı
taşları !!{signed formula}s, yani bir doğruluk değeri ``işareti'' ($\True$
ya da~$\False$) taşıyan !!{sentence}s{}dir. Bu işaretli !!{formula}s (aşağıya
doğru büyüyen) bir ağaçta düzenlenir.

\begin{defn}
  Bir \emph{!!{signed formula}}, bir doğruluk değeri ile !!a{sentence}'den
  oluşan bir ikilidir; başka bir deyişle şu iki biçimden birindedir:
  \[
  \sFmla{\True}{!A} \text{ veya } \sFmla{\False}{!A}.
  \]
\end{defn}

Sezgisel olarak $\sFmla{\True}{!A}$ ifadesini ``$!A$ doğru olabilir'',
$\sFmla{\False}{!A}$ ifadesiniyse ``$!A$ yanlış olabilir'' (belirli bir
\iftag{FOL}{!!{structure}}{değerleme} içinde) diye okuyabiliriz.

Ağaçtaki her !!{signed formula} ya bir \emph{varsayım}dır (bunlar ağacın en
üstünde sıralanır) ya da üstündeki !!a{signed formula}'dan çeşitli çıkarım
kurallarından biriyle elde edilir. Önceki ifadenin olası her !!{main
operator} için iki kural vardır; burada söz konusu ifade önceki
!!{formula}'dır: işaretin~$\True$ olduğu durum için bir kural ve
işaretin~$\False$ olduğu durum için bir kural. Bazı kurallar
ağacın dallanmasına izin verirken bazıları yalnızca dala !!{signed formula}s
ekler. Bir kural, hemen önce gelen !!{signed formula}'ya değil, kökten
kuralın uygulandığı yere uzanan daldaki herhangi bir !!{signed formula}'ya
uygulanabilir (ve çoğu zaman uygulanmalıdır).

Bir dal hem $\sFmla{\True}{!A}$ hem de $\sFmla{\False}{!A}$ içerdiğinde
\emph{kapalı}dır. Her dalı kapalı olan !!{tableau}, kapalıdır. Sezgisel
yorum uyarınca her dal ortaklaşa mümkün bir durumu betimler; ancak
$\sFmla{\True}{!A}$ ile $\sFmla{\False}{!A}$ ortaklaşa mümkün değildir.
Başka bir deyişle, bir dal kapalıysa betimlediği olasılık elenmiştir.
Özellikle bu, kapalı bir !!{tableau}'nun $\sFmla{\True}{!A}$ biçimindeki her
varsayımı aynı anda doğru ve~$\sFmla{\False}{!A}$ biçimindeki her varsayımı
aynı anda yanlış kılmaya yönelik bütün olasılıkları elediği anlamına gelir.

$!A$ \emph{için kapalı bir !!{tableau}}, kökü~$\sFmla{\False}{!A}$ olan
kapalı bir !!{tableau}'dur. Böyle kapalı bir !!{tableau} varsa $!A$'nın
yanlış olmasına yönelik bütün olasılıklar elenmiştir; yani $!A$ her
\iftag{FOL}{!!{structure}}{değerleme} içinde doğru olmalıdır.

